\section{Failures and corrections}
\label{sec:failures}

The protocol did not prevent mistakes. It made them visible, and it fixed in
advance what each one would cost. This section records the ones that changed
a result or a claim.

\subsection{Leaks that realism gates could not see}

The event simulator passed its preregistered realism gates, per-period event
rates and score effects, with two leaks present. The first was an orientation
flag equal to the team that owned the \emph{next} event, which is the team
component of the prediction target; it agreed with the target's team in every
held-out step. The second was worse. The model conditioned on the on-ice
players at the next event, and because line changes happen at stoppages, a
change in on-ice personnel announces the faceoff that caused it: when the
next event's on-ice set lost one skater, the next event was a faceoff with
probability 0.986 against a base rate of 0.199. About 3,000 of 62,094 held-out
steps were given a near-deterministic answer.

Neither gate could catch either leak, because both scored the model on the
inputs it had been trained with, and a leaked input improves realism along
with everything else. The first leak surfaced only when a downstream consumer
had to supply the flag itself and produced a home scoring rate of 0.39 goals
per 60 minutes against 3.30 for the away side; the second surfaced in the audit
that followed. The fix conditions on the current on-ice set and fixes the
orientation, which lowers the actor head's ceiling honestly: the next actor is
on the ice for 59.6\% of steps, against the 82.7\% the leaked version enjoyed.
All earlier results for the simulator were voided and every seed was
retrained. The lesson generalised into the standing causality audit of
Section~\ref{sec:protocol}: leakage needs either a consumer that must supply
every input itself or a direct test of causality; realism is not evidence of
legitimacy.

\subsection{A measurement drift}

The shot-quality model's calibration gate failed three times, and each failure
was diagnosed on training-side seasons before the next correction was
declared. First, the calibrator had been fitted on cross-validation folds and
applied to a model trained on all the data, which was better calibrated than
the folds, so calibration made it worse (maximum decile deviation 0.0088
against 0.0061 uncalibrated on a development vantage). Second, the recording
change of Section~\ref{sec:data} moved the meaning of a shot location from 2023
on; a calibrator frozen at the end of the previous season is exactly one
season behind a moving target. The declared remedy refits the isotonic map
within the season on trailing out-of-sample shots. Third, per-rink recording
offsets were added as walk-forward covariates. The pooled maximum decile
deviation fell from 0.0222 to 0.0124 and then to between 0.0060 and 0.0074
across the later declared stages, but never reached the 0.005 bar. The failure
is recorded against the transition itself, which no walk-forward method can
see before it happens.

The same drift reached \neurhlG{}. Trained through 2024, it learned the older ratio
of goals to expected goals and reads 2025--26-era expected goals as inputs, so
a dry run on opening night projected 3.36 regulation goals per team-game
against a league level near 3.0. Because the win-probability stack is fitted on
the engine's unscaled outputs and the box-score simulation conditions every
score on the stacked outcome, the win probabilities are unaffected; only the
stat sheet's scoring level is. A declared amendment, committed before any
2026--27 game, scales the stat sheet's goal means by a factor fixed before the
season and updated walk-forward (Section~\ref{sec:live}).

\subsection{Protocol corrections}

\paragraph{A rule changed with the result in view.}
The first plan's fourth amendment required every headline result to be
reported both with and without the anomalous seasons. A later plan replaced
that rule with a simpler one, train but never score 2013 and 2021, and called
the replacement a verbatim carry-over. It was not: the replacement was written
after tune-window results were visible, and the commit that introduced it also
recorded \neurhlH{}'s tune-window pass ($p = 0.0075$). Between the committed tune
failure ($p = 0.120$) and that pass, three defensible things had changed on the
same window: the scored seasons, two corrections to Layer~1, and the scoring
rule. Together they were decisions taken with the result in view. A later
amendment reclassified the tune-window result as exploratory, required it to
be reported under every scored set ($p$ from 0.0075 to 0.019), and declared the
one-shot confirmation reported in Section~\ref{sec:res-h}.

\paragraph{A comparison across different seasons.}
An early summary stated that the season layer beat both house benchmarks. The
benchmark figure came from 2018--2026 and the \neurhl{} figure from 2011--2017.
On matched seasons the Elo baseline is better (Section~\ref{sec:res-season});
the claim is withdrawn.

\paragraph{Reliability is not validity.}
The first RAPM gate required split-half reliability above that of raw on-ice
shot share, and RAPM failed it (offence 0.656 against 0.789). The gate was
wrong: raw on-ice rates are more repeatable because they are contaminated, as
a player keeps the same linemates, zone starts and competition in both halves,
and RAPM removes exactly those terms. Ranking rating systems by reliability
rewards confounding. The gate was retired and replaced by predictive transfer
to a new team, on which RAPM tied a one-line team-demeaned baseline.

\paragraph{Gates implemented short of their declaration.}
While writing this paper we found that the script implementing \neurhlG{}'s gates
had omitted three declared components: the randomised-PIT coverage criterion
of the calibration gate, the comparison against the summed player-game layer,
and CRPS. The calibration gate had been judged on its slope and overtime share
alone and recorded as a pass. Computed afterwards from the same saved
predictions, with nothing refitted, the team-box-score comparisons pass but the
coverage for team shots is 0.860 against a nominal 0.80 and a tolerance of 0.03,
so the calibration gate fails as declared. The correction was appended to the
plan as a dated amendment. It changes no decision, because the pre-gate had
already failed on its \neurhlH{} condition, but it changes a public claim, and
Table~\ref{tab:gates} reports the corrected record.
